\documentclass[11pt]{article}

\usepackage[margin=1in]{geometry}
\usepackage{booktabs}
\usepackage{amsmath}
\usepackage{hyperref}

\title{Matched-Frequency Synthetic Relations: Structured Affine Relation vs Random Memorized Relation}
\author{}
\date{}

\begin{document}
\maketitle

\section{Question}

We test whether two relations with identical frequency can be learned by different
mechanisms.  Relation \(R_0\) is supported by latent attribute structure. Relation
\(R_1\) is a random permutation over the same subjects and objects. Both relations
have the same number of distinct trained triples and the same repetition count.

\section{Configuration}

\begin{table}[h]
\centering
\begin{tabular}{ll}
\toprule
Quantity & Value \\
\midrule
Subjects / objects & 192 / 192 \\
Attributes & \(f \in \{0,\ldots,31\}\), \(g \in \{0,\ldots,31\}\) \\
Structured shift & \((\Delta_f,\Delta_g) = (3,5)\) \\
Relations & \(R_0\) structured, \(R_1\) random \\
Trained triples per relation & \(k = 128\) \\
Held-out subjects & 64 \\
Repeats per trained triple & 128 for \(R_0\), 128 for \(R_1\) \\
Attribute sentence repeats & 4 \\
Model & 1-layer causal transformer \\
Hidden dimension / heads & 64 / 4 \\
Context length & 8 \\
Embedding tying & false \\
Learning rate & \(10^{-3}\) \\
Object loss weight & 4.0 \\
\bottomrule
\end{tabular}
\caption{Final matched-frequency configuration.}
\end{table}

The run is configured in
\texttt{experiments/configs/runs/synthetic\_mixed\_relations\_category\_k128\_r128.toml}.
Training is ordinary end-to-end gradient descent. No components are frozen and no
targeted component updates are used.

\section{Data Generation}

Each subject token \(a_i\) is assigned two latent attributes:
\[
  a_i \mapsto (f_i,g_i).
\]
The structured target is a job-category object determined by the subject's
\(f\)-attribute:
\[
  R_0(a_i) = b_{f_i}.
\]
Thus many subjects can share the same structured target, analogous to many applicants
with the same skill bucket matching the same job category.

For the category objects \(b_j\), \(j < 32\), object attributes are shifted:
\[
  b_j \mapsto (j + 3 \bmod 32,\; 0).
\]
Remaining object tokens are distractor objects used by the random relation.

The model observes attribute sentences for all subjects and objects:
\[
  x\; a_i\; f_i\; g_i\; x,
\]
\[
  x\; b_i\; f'_i\; g'_i\; x.
\]

The structured relation is therefore supported by the latent skill/category
attribute \(f_i\).

The random relation is
\[
  R_1(a_i) = b_{\pi(i)},
\]
where \(\pi\) is a seed-dependent random permutation. This relation is not supported
by the attributes.

The relation training sentences are
\[
  x\; a_i\; r_0\; b_{f_i}\; x,
\]
\[
  x\; a_i\; r_1\; b_{\pi(i)}\; x.
\]

A seed-dependent shuffle selects 128 subjects for relation training and 64 subjects
for held-out evaluation. Held-out subjects still appear in attribute sentences, but
never appear in relation sentences.

\section{Affine Test}

Affine structure is evaluated with ordinary least-squares linear regression. For each
relation, we fit
\[
  \hat{o} = W s + b
\]
from subject token embeddings \(s\) to the target object token embeddings \(o\), using
the 128 trained relation subjects. We then evaluate on the 64 held-out subjects.

The main metric is affine nearest-neighbor accuracy: the predicted vector \(\hat{o}\)
is compared against all object embeddings, and the prediction is correct if the
nearest object is the held-out target.

\section{Main Result}

\begin{table}[h]
\centering
\begin{tabular}{lccccc}
\toprule
Relation & Train top-1 & Held-out top-1 & Held-out top-10 & OLS affine NN & OLS affine cosine \\
\midrule
\(R_0\) structured & 1.000 & 0.875 & 0.953 & 0.969 & 0.910 \\
\(R_1\) random & 1.000 & 0.000 & 0.000 & 0.016 & 0.214 \\
\bottomrule
\end{tabular}
\caption{Seed 19. Both relations are learned on trained triples; only \(R_0\) supports
linear-regression affine generalization.}
\end{table}

\section{Replication}

\begin{table}[h]
\centering
\begin{tabular}{rcccccc}
\toprule
Seed & \(R_0\) train & \(R_0\) held-out & \(R_0\) affine NN &
\(R_1\) train & \(R_1\) held-out & \(R_1\) affine NN \\
\midrule
31 & 1.000 & 0.641 & 0.906 & 1.000 & 0.000 & 0.000 \\
47 & 1.000 & 0.859 & 0.812 & 1.000 & 0.000 & 0.000 \\
59 & 1.000 & 0.797 & 0.938 & 1.000 & 0.000 & 0.000 \\
\bottomrule
\end{tabular}
\caption{Replication seeds preserve the same pattern.}
\end{table}

\section{Interpretation}

The frequency structure is matched exactly: both relations have 128 distinct trained
triples and each trained triple is repeated 128 times. Both relations reach perfect
top-1 recall on trained triples. \(R_0\) also generalizes behaviorally to held-out
subjects, while \(R_1\) does not. Ordinary linear regression recovers the held-out
structured mapping for \(R_0\), while it fails for the random mapping \(R_1\).

Thus, under matched frequency, the attribute-supported relation is represented with
substantial affine structure, whereas the random relation is memorized without affine
generalization.

\section{Reproduction}

\begin{verbatim}
cd experiments
env UV_CACHE_DIR=/tmp/uv-cache MPLCONFIGDIR=/tmp/mpl-cache \
  uv run factual-recall synthetic-mixed-relations \
  --config configs/runs/synthetic_mixed_relations_category_k128_r128.toml \
  --device cpu --quiet
\end{verbatim}

\end{document}
